\section{Proposed PhysFormer Framework}
\label{sec:method}

\subsection{Mathematical Formulation}
Throughout this paper, we adopt the following notation: $P_{\text{load}}$, $P_{\text{pv}}$, and $P_{\text{wind}}$ denote load, PV, and wind power [MW]; $G$, $T$, and $v$ represent irradiance [W/m$^2$], temperature [$^\circ$C], and wind speed [m/s]; $\eta_{\text{pv}}$ and $\beta_T$ are the PV conversion efficiency and temperature coefficient; $v_{\text{ci}}$, $v_r$, and $v_{\text{co}}$ are the wind turbine cut-in, rated, and cut-out speeds; $a_{\text{pv}}$ and $a_{\text{wind}}$ are physics activity masks $\in[0,1]$; $g_{\text{pv}}$ and $g_{\text{wind}}$ are PGCC gates; $\alpha$ is the curriculum factor; $S=672$ and $P=96$ are the sequence and prediction lengths; and $D=512$ is the model dimension.

In the context of Virtual Power Plant (VPP) forecasting, the objective is to predict the future power outputs of three distinct sub-components: base load ($P_{\text{load}}$), photovoltaic generation ($P_{\text{pv}}$), and wind generation ($P_{\text{wind}}$). The overall net power demand of the VPP is defined as $P_{\text{net}} = P_{\text{load}} - P_{\text{pv}} - P_{\text{wind}}$. Let $\mathcal{X}_{\text{stat}} \in \mathbb{R}^{B \times S \times 3}$ denote the historical observation sequence for the three components over a look-back window of length $S$, where $B$ is the batch size. Concurrently, the corresponding historical weather conditions---comprising temperature, irradiance, and wind speed---are represented by $\mathcal{X}_{\text{weather\_hist}} \in \mathbb{R}^{B \times S \times 3}$. The model is also provided with future weather forecasts for the prediction horizon $P$, denoted by $\mathcal{X}_{\text{weather\_future}} \in \mathbb{R}^{B \times P \times 3}$. Given these inputs, the forecasting model aims to generate point predictions for the output $Y \in \mathbb{R}^{B \times P \times 3}$, providing the respective future trajectories for the load, PV, and wind components.

\subsection{Overall Architecture}

\begin{figure}[htbp]
\centering
\includegraphics[width=\columnwidth]{IEEE_Fig1_Architecture.png}
\caption{The overall architecture of the proposed PhysFormer. It features a dual-stream design: a statistical stream encoded by a Transformer and an explicit physical mapping stream. These streams are fused through the Physics-Guided Gating and Alignment (PGCC) module. Future weather conditions are injected using a Gated Linear Unit (GLU) to provide explicit foresight. Finally, theory-anchored residual heads refine the predictions under physics-derived activity gates to reduce non-physical outputs.}
\label{fig:overall_architecture}
\end{figure}

The PhysFormer is constructed as a dual-stream Transformer with PGCC and theory-anchored activity gating. The processing pipeline follows seven steps: (1) Historical observation $\mathcal{X}_{\text{stat}}$ and weather $\mathcal{X}_{\text{weather\_hist}}$ are concatenated and encoded by a Transformer encoder ($d_{\text{model}}=512$, $d_{\text{ff}}=1024$, $8$ heads, $3$ layers). (2) $\mathcal{X}_{\text{weather\_hist}}$ is processed through ExplicitPhysicalMapping for physical baselines. (3) PGCC fuses statistical and physics representations. (4) FlattenHead projects from historical dimension $S$ to prediction horizon $P$. (5) Future weather $\mathcal{X}_{\text{weather\_future}}$ is injected via a GLU with residual logic. (6) Separate output heads predict per-component residual corrections. (7) Theory-Anchored Output applies physics activity masks to yield final predictions.

\subsection{Explicit Physical Mapping}
The core of the physics injection lies in the ExplicitPhysicalMapping layer, which maps relevant weather variables to theoretical power baselines via domain-specific equations.

\subsubsection{PV Model}
The theoretical PV generation is driven by irradiance $G$ and degraded by cell temperature, modeled as:
\begin{equation}
    P_{\text{pv}} = G \cdot \eta_{\text{pv}} \cdot \text{clamp}(1 - \beta_T \cdot (T - 25), 0, 1.5)
\label{eq:pv_model}
\end{equation}
where $T$ is the temperature, $\eta_{\text{pv}} = 0.33 \cdot \sigma(\theta_{\text{eff}})$ is the learnable conversion efficiency bounded by the Shockley-Queisser limit, and $\beta_T = \text{softplus}(\theta_{\text{temp\_coef}}) \times 0.01 \in [0.003, 0.005]/^\circ\text{C}$ is the temperature coefficient.

\subsubsection{Wind Model}
Wind power limits depend heavily on operating thresholds: the cut-in speed $v_{\text{ci}}$, rated speed $v_{\text{r}}$, and cut-out speed $v_{\text{co}}$. To preserve strict threshold ordering, we utilize a cumulative delta structure where $v_{\text{ci}} = \text{softplus}(\theta_{\text{ci}})$, $v_{\text{r}} = v_{\text{ci}} + \text{softplus}(\theta_{\text{r\_delta}})$, and $v_{\text{co}} = v_{\text{r}} + \text{softplus}(\theta_{\text{co\_delta}})$. The power output uses an idealized soft-curve approximation modified by operational compliance:
\begin{equation}
    P_{\text{wind}} = \text{scale}_{\text{wind}} \cdot \sigma(5 \cdot (w_{\text{norm}} - 0.5)) \cdot I_{\text{running}}
\label{eq:wind_ideal}
\end{equation}
where $w_{\text{norm}}$ is the normalized speed, $scale_{\text{wind}}$ is a learnable amplitude factor, and $I_{\text{running}}$ is 1 within $[v_{\text{ci}}, v_{\text{co}}]$ and 0 otherwise. To maintain differentiability, PhysFormer replaces $I_{\text{running}}$ with a continuous physics activity mask $a_{\text{wind}}$ (detailed in Section~\ref{subsec:activity_gating}):
\begin{equation}
    P_{\text{wind,actual}} = \text{scale}_{\text{wind}} \cdot \sigma(5 \cdot (w_{\text{norm}} - 0.5)) \cdot a_{\text{wind}}
\label{eq:wind_actual}
\end{equation}
This decoupling bridges rigid physical thresholds with end-to-end backpropagation without Straight-Through Estimators. The cumulative delta structure ensures overlapping bounds never collapse during gradient updates.

\subsubsection{Load Model}
We abstract the power load variations triggered by human cooling and heating activities through a quadratic thermal comfort model:
\begin{equation}
    P_{\text{load}} = P_{\text{base}} + \text{softplus}(\gamma) \cdot (T - T_c)^2
\label{eq:load_model}
\end{equation}
where $T_c$ is an optimizable thermal comfort baseline initialized at $20^\circ\text{C}$. While the self-attention mechanism of the Transformer is well suited to extracting the diurnal and weekly periodicities that characterize human social activities, the Explicit Physical Mapping stream is deliberately restricted to isolating the non-linear, temperature-driven volatility of the electrical load. Therefore, we conceptualize the thermal-induced variance through this quadratic comfort model, explicitly decoupling meteorological sensitivity from the statistically captured chronological baselines.

\begin{figure*}[t]
\centering
\includegraphics[width=0.85\textwidth]{IEEE_Fig2_PhysicalMapping.pdf}
\caption{Learned explicit physical mapping functions after convergence. (a) PV temperature degradation under varying irradiance levels ($\beta_T = 0.0041/^\circ$C). (b) Wind turbine power curve with smooth sigmoid transition. (c) Load thermal comfort demand ($T_c \approx 20^\circ$C). Parameters are listed in Table~\ref{tab:physical_params}.}
\label{fig:physical_mapping}
\end{figure*}

\subsubsection{Physics Activity Gating}
\label{subsec:activity_gating}
To prevent non-physical artifacts, we propose an explicit Activity Gating scheme. For PV, night-time irrelevance: $a_{\text{pv}} = \tanh(10 \cdot P_{\text{pv\_theory}})$. For wind, speed-dependent activity with differentiable soft limits: $a_{\text{wind}} = \tanh(10 \cdot \text{softplus}(v - v_{\text{ci}}) \cdot \text{softplus}(v_{\text{co}} - v))$. Within the BPAR framework, the residual mechanism automatically compensates for systematic dampening effects near thresholds. The load takes an unconditional $a_{\text{load}} = 1$. These physics-derived masks maintain continuous gradient flow while bridging active operating states and null-generation periods.

\subsection{Physics-Guided Gating and Alignment (PGCC)}
The Physics-Guided Gating and Alignment (PGCC) module coordinates the interaction between the statistical stream and the physically encoded representations. In this work, PGCC is used to guide weather-response behavior and improve alignment with known physical cues rather than to infer causality from data alone.

\subsubsection{Physics-Anchored Illumination Prior}
We first establish a data-guided illumination prior to suppress spurious activation patterns. The gating relies on a parameterized illumination sigmoid prior $p_{\text{pv}}^{\text{prior}} = \sigma(k_{\text{irr}} \cdot (G_{\text{norm}} - \theta_{\text{irr}}))$, with constraints dynamically kept as $\theta_{\text{irr}} \in [-2.0, 0.5]$ and bounded slope $k_{\text{irr}} \in [1, 100]$. This prior helps distinguish day and night intervals and reduces the influence of noisy activations in the data embeddings.

\subsubsection{Volatility-Aware Soft Gate}
To dynamically scale network trust, the system embeds sequence-level sample fluctuations extracted as $\upsilon = \frac{1}{C}\sum_{c}\text{Var}_{t}(\mathcal{X}_{\text{weather},c})$. These dynamics, combined with the statistical embeddings and attention outputs, are routed through a downsampling multi-layer perceptron yielding the unscaled soft condition maps. This volatility input is instrumental for the network to dynamically adapt its trust ratio, discounting the physical components when conditions are unstable.

\subsubsection{Source-Differentiated Gate Bounds}
The coupling implements separate activation ranges to accommodate generation source characteristics:
\begin{gather}
    \text{gate}_{\text{pv}} = \text{smooth} \left( p_{\text{pv}}^{\text{prior}} \cdot (0.5 + g_{\text{soft\_pv}}) \right) \in [0, 1.5] \label{eq:gate_pv} \\
    \text{gate}_{\text{load}} = \text{smooth}(g_{\text{soft\_load}}) \in [0, 1.0] \label{eq:gate_load}
\end{gather}
The asymmetric upper bound of $1.5$ for $\text{gate}_{\text{pv}}$ provides \textit{engineering headroom} for the neural residual within the BPAR additive structure. During meteorological transients (e.g., Cloud Edge Effect with $1.2\times$--$1.5\times$ over-irradiance), this ceiling preserves sufficient dynamic range for the residual stream to recover transient peaks. Through explicit gate response regularization, the learned $\text{gate}_{\text{pv}}$ reaches a Pearson correlation of $r=0.8306$ with exogenous irradiance (Fig.~\ref{fig:causal_gate}), indicating that the gate remains aligned with the corresponding weather cue while retaining predictive flexibility.

\begin{figure}[htbp]
\centering
\includegraphics[width=\columnwidth]{IEEE_Fig3_CausalGate.pdf}
\caption{Gate-irradiance alignment in PGCC. Top: time-series overlay of learned PV gate and solar irradiance over representative test samples. Bottom: scatter regression on a visualization subset; the global Pearson correlation across the full test set is $r=0.8306$ ($p < 0.001$).}
\label{fig:causal_gate}
\end{figure}

\subsubsection{Physics-to-Data Curriculum Transition}
We condition the alignment phase via a curriculum factor $\alpha$. As the fusion queries retrieve features, the module integrates inputs dynamically scaled as:
\begin{equation}
    q = \alpha \cdot F_{\text{physics}} + (1 - \alpha) \cdot g_{\text{soft}} \cdot F_{\text{stat}}
\label{eq:curriculum_query}
\end{equation}
During the initial training phases ($\alpha \to 1$), the cross-attention queries rely more heavily on the structured physical representations. As training progresses ($\alpha \to 0$), the queries gradually transition toward the data-learned residual states gated by $g_{\text{soft}}$, balancing physics-guided initialization with empirical statistical adaptation.

\subsubsection{Temporal Smoothing}
Due to internal frame dependencies, local sequence discontinuities often plague raw gating parameters. To counteract this, all temporal gating variables pass through a grouped one-dimensional depthwise convolution (kernel size $3$, groups=$d_{\text{model}}$). This filter structurally suppresses high-frequency chaotic fluctuation within the gate assignment layer while preserving macro-level shift indicators.

\subsection{Exogenous Weather Integration}
After the FlattenHead projects the fused representations into the explicit future context, the framework must systematically integrate upcoming weather patterns at the given horizon. Because the prediction horizon's meteorological conditions are unavailable in the historical context alone, providing this explicit foresight is essential. The future physical latent space $F_{\text{weather\_future}}$ is concatenated with the neural history $F_{\text{history}}$, and aggregated via a Gated Linear Unit ($\text{gate} = \sigma(\text{linear}(F_{\text{cat}}))$). The bounded residual output is produced as $F_{\text{future}} = \text{gate} \cdot \text{proj}(F_{\text{cat}}) + F_{\text{history}}$. This explicitly bounded forward visibility ensures the residual heads are properly conditioned to react to impending deterministic disruptions, such as dawn transitions or sudden wind cut-outs.

\subsection{Bounded Physical Act-Residual (BPAR) Mechanism} \label{subsec:bpar}
To systematically mitigate non-physical anomalies while preserving the continuous gradient flow, we propose the Bounded Physical Act-Residual (BPAR) mechanism. Let $\mu_x$ and $\sigma_x$ denote the global mean and standard deviation of component $x$. We define the normalized physical zero-point as $P_{x,\text{zero}} = (0 - \mu_x) / \sigma_x$. BPAR reconstructs the predictive manifold by anchoring the neural residual $r_x$ onto the theory baseline within a differentiable Softplus manifold:
\begin{equation}
    \hat{P}_{x,\text{raw}} = \text{Softplus}\Big((P_{x,\text{theory}} + r_x) - P_{x,\text{zero}}\Big) + P_{x,\text{zero}}
\label{eq:bpar}
\end{equation}
By the fundamental property of Softplus ($\text{Softplus}(z) > 0, \forall z$), this parameterization keeps $\hat{P}_{x,\text{raw}} > P_{x,\text{zero}}$ in the normalized domain. Under the adopted normalization and anti-normalization setting, $\hat{P}_{\text{phys}} = \hat{P}_{x,\text{raw}} \cdot \sigma_x + \mu_x > 0$~MW for $\sigma_x > 0$, which gives the model a structurally non-negative output form. We maintain a numerical floor $\epsilon=10^{-5}$ for degenerate variance states (e.g., night-time PV), supporting consistent boundary compliance in practice.
Subsequently, the prediction is fused with the physics activity mask $a_x \in [0, 1]$:
\begin{equation}
    \hat{P}_{x,\text{final}} = a_x \cdot \hat{P}_{x,\text{raw}} + (1 - a_x) \cdot P_{x,\text{zero}}
\label{eq:activity_mask}
\end{equation}
Whenever $a_x \to 0$ (driven by meteorological priors), BPAR smoothly drives output toward the physical zero baseline and reduces boundary violations. Compared with unconstrained predictors that may cross the zero-power boundary, BPAR favors physically consistent outputs at the cost of small local corrections near the boundary.

\subsection{Physics-Constrained Curriculum Training}
\label{subsec:training}

Training deep neural networks with rigid physical constraints poses significant optimization challenges due to gradient conflicts between data-driven and physics objectives. We design a Physics-Constrained Curriculum Training strategy that decouples metric spaces, enforces dynamic amplitude balancing, and maps learning trajectories onto a phased schedule.

\subsubsection{Loss Function Design}
The primary regression loss (MSE) operates in normalized space for stable gradients, while all physics constraints ($L_{\text{phys}}$) are evaluated after denormalization into the physical domain (MW). The physical objectives comprise six penalties (Table~\ref{tab:losses}): soft constraints ($L_{\text{net}}$, $L_{\text{energy}}$, $L_{\text{deriv}}$, $L_{\text{dir}}$) guide structural topology, while hard constraints ($L_{\text{bvr}}$, $L_{\text{rvr}}$) use aggressive quadratic formulations for safety-critical boundaries.

\begin{table}[htbp]
\caption{Components of Physics-Guided Loss Function}
\label{tab:losses}
\centering
\footnotesize
\setlength{\tabcolsep}{3pt}
\begin{tabularx}{\columnwidth}{@{}llc>{\raggedright\arraybackslash}X@{}}
\toprule
\textbf{Comp.} & \textbf{Wt.} & \textbf{Formula} & \textbf{Motivation} \\ \midrule
$L_{\text{net}}$ (Soft) & 1.0 & $|\sum \hat{y}_{\text{comp}} - y_{\text{net}}|$ & Ensures models sum to net demand. \\
$L_{\text{energy}}$ (Soft) & 0.5 & $|\mathbb{E}_t[\hat{y}] - \mathbb{E}_t[y]|$ & Enforces macro energy conservation. \\
$L_{\text{deriv}}$ (Soft) & 0.3 & $|\Delta \hat{y} - \Delta y|$ & Captures ramp rate magnitudes. \\
$L_{\text{dir}}$ (Soft) & 0.1 & $\mathbb{E} [1 {-} \text{sim}_{\text{cos}}]$ & Preserves accurate trend switching. \\ \midrule
$L_{\text{bvr}}$ (Hard) & 2.0 & $\mathbb{E} [(\max(0, -\hat{y}_{\text{real}}))^2]$ & Penalizes negative power strictly. \\
$L_{\text{rvr}}$ (Hard) & 2.0 & $\mathbb{E} [(\max(0, |\Delta\hat{y}| {-} \rho_k))^2]$ & Penalizes sudden jumps. \\ \bottomrule
\end{tabularx}
\end{table}

We propose an adaptive amplitude balancing factor $\gamma_{\text{ema}}$ that dynamically tracks the ratio between current forecasting inaccuracy and physical compliance:
\begin{equation}
    \gamma_{\text{ema}} \leftarrow 0.9 \cdot \gamma_{\text{ema}} + 0.1 \cdot \left( \frac{L_{\text{mae\_phys}}}{L_{\text{phys\_ref}}} \right)
\label{eq:gamma_ema}
\end{equation}
where $L_{\text{phys\_ref}}$ is the total unweighted physical loss in MW. This yields a dimensionless multiplier, defaulting to $1.0$ for the first 50 batches, halting updates when curriculum activation $< 0.05$, and clamped to $[0.1, 10.0]$. The combined objective becomes:
\begin{equation}
    L_{\text{total}} = L_{\text{mse}}(\text{normalized}) + \gamma_{\text{ema}} \cdot L_{\text{phys\_weighted}}(\text{MW})
\label{eq:total_loss}
\end{equation}

Note: while BPAR provides a structural output floor, $L_{\text{bvr}}$ remains essential during training to provide gradient signals that prevent latent representations from drifting into non-physical regimes.

\subsubsection{Four-Phase Curriculum}
To bridge the parameter convergence gap between the arbitrary Gaussian distributions at neural initialization and the highly structured physical reality, we inject a Four-Phase constraint schedule, summarized in Table~\ref{tab:curriculum}.

\begin{table}[htbp]
\caption{Four-Phase Physics Curriculum and Layer-wise Learning Rates}
\label{tab:curriculum}
\centering
\footnotesize
\setlength{\tabcolsep}{2pt}
\begin{tabularx}{\columnwidth}{@{}lc>{\raggedright\arraybackslash}p{1.8cm}>{\raggedright\arraybackslash}X@{}}
\toprule
\textbf{Phase} & \textbf{Ep.} & \textbf{Active Const.} & \textbf{Notes} \\ \midrule
I (Init.) & 0--4 & None & Physics frozen; base mapping. \\
II (Safety) & 5--14 & $L_{\text{net}}, L_{\text{bvr}}$, $L_{\text{rvr}}, L_{\text{dir}}$ & Safety bounds ramped. \\
III (Integ.) & 15--29 & Phase II + $L_{\text{energy}}, L_{\text{deriv}}$ & Macro conservation. \\
IV (Satur.) & 30--99 & All & Full regularized training. \\ \midrule
\multicolumn{4}{@{}l}{\textbf{Layer-wise Learning Rates}} \\ \midrule
\multicolumn{2}{@{}l}{Physics (ExplicitPhysicalMapping)} & $10^{-5}$ & Preserves priors. \\
\multicolumn{2}{@{}l}{Gating (PGCC)} & $5{\times}10^{-5}$ & Stable alignment. \\
\multicolumn{2}{@{}l}{Statistical (all other)} & $10^{-4}$ & Fast convergence. \\
\bottomrule
\end{tabularx}
\end{table}

This curriculum progressively introduces physical-compliance objectives. Phase I isolates statistical representations with frozen physics parameters. Phase II prioritizes safety boundaries. Phase III introduces macroscopic features, and Phase IV activates the full regularized objective. Each constraint intensity is transitioned via a bounded sigmoid ramp $\sigma(12 \cdot (\text{progress} - 0.5))$. The PGCC alignment coefficient follows $\alpha = 1.0 - w_{\text{curr\_net}}$, coupling the internal attention fusion with the external physical guidance term.

\subsubsection{Layer-wise Learning Rates}
To sustain learned feature fidelity without abruptly disrupting the theory-anchored baselines pre-defined by human engineering inputs, we decouple the backpropagation update amplitudes into three isolated AdamW groupings (Table~\ref{tab:curriculum}, bottom).

The baseline parameters proceed at a default learning rate of $1 \times 10^{-4}$. After their suspension ends in epoch 5, the core generation physics operators unfreeze at a drastically suppressed $1 \times 10^{-5}$ ($0.1\times$ the baseline rate). This strict limitation permits minor optimization fine-calibration over coefficients like optimal irradiance generation thresholds while halting gross semantic reconstruction. Correspondingly, the intermediate gating features proceed at $5 \times 10^{-5}$, which yields smooth adjustment profiles.

\subsubsection{Regularization and Early Stopping}
Three regularization pathways reinforce generalizability. First, a Physics Prior Regularizer evaluates learned physical properties against nameplate values via MSE (weight 0.1). Second, a Gate Response Regularization anchors the temporal PV gating variable to external irradiance (weight 0.05), with dynamic masking that shuts down penalization when prior variance $\sigma^2_{\text{prior}} < 10^{-3}$ (bypassing night-time noise).

\textbf{Role of PGCC:} In the present framework, PGCC mainly improves gate alignment and interpretability rather than raw forecasting accuracy. Because the final prediction $\hat{P}_x = \text{Softplus}(P_{\text{theory}} + r_x)$ is still produced through the residual pathway, the supervised alignment term can shape weather-response behavior without preventing the residual heads from minimizing prediction error.

Gradient clipping uses a dynamic bound of $0.5$ during critical curriculum transitions ($0.1 < w_{\text{curr\_net}} < 0.9$), releasing to $1.0$ otherwise. Early stopping uses channel-equalizing $NRMSE_{\text{avg}}$ with patience of 15 epochs.
